\subsection{Two executed contour certificates}
\label{sec:executed-contour}
The error bounds above can be evaluated using the defining zeta
function and whole-edge complex enclosures. We give two such
certificates, at positive ordinate heights $20$ and $45$. The
coordinate is still $z=2\gamma$.

We first give a uniform elementary bound for the heat perturbation.
For $u\ge0$ the kernel in Section~\ref{sec:fd-positive-surgery} satisfies
\[
 0<\Phi(u)
 \le 2\pi^2 e^{9u-\pi e^{4u}}
       \sum_{n\ge1}n^4e^{-\pi(n^2-1)e^{4u}}
 <40e^{9u-3e^{4u}}.
\]
Indeed, the sum is less than $2$. Its terms with $n\ge2$ are
bounded by $n^4e^{-3(n^2-1)}$; the first is less than $1/32$,
and successive ratios are less than
$(3/2)^4e^{-15}<1/1000$. Since
$e^{4u}\ge1+4u+8u^2$, it follows that
\[
 u^j e^{u^2+u/4}\Phi(u)
 <40e^{-3}u^j e^{-23u^2-11u/4},\qquad 0\le j\le3.
\]
Each integral on the right is less than $1$: drop its last
exponential and use, respectively,
\[
 \int_0^\infty u^j e^{-23u^2}\,du
 =\frac{\sqrt\pi}{2\sqrt{23}},\quad
 \frac1{46},\quad
 \frac{\sqrt\pi}{4\,23^{3/2}},\quad
 \frac1{2\,23^2},
\]
together with $40e^{-3}<5/2$, $\sqrt\pi<2$ and
$\sqrt{23}>4$. Differentiating the cosine integral therefore proves
\begin{equation}\label{eq:uniform-contour-heat}
 |H_t(z)-H_0(z)|\le |t|,\qquad |\partial_zH_t(z)|<1
 \quad (|t|\le1,\ |\Im z|\le1/4).
\end{equation}

For $T=20,45$, let $U_T$ be the rectangle
$0<\Re z<2T$, $|\Im z|<1/4$. The certified data are
\[
\begin{array}{c|c|c|c|c}
 T&\min_{\partial U_T}|H_0|\text{ exceeds}&R_T&A_T\text{ is below}
       &S_0(0)\\ \hline
 20&10^{-6}&10^{-8}&14\cdot10^6&1\\
 45&10^{-14}&10^{-16}&3\cdot10^{15}&8
\end{array}
\]
Here $A_T$ has exactly the definition in the preceding subsection.
The boundary length is $4T+1$, so
\[
 A_T\le\frac{4T+1}{2\pi(\eta_T-R_T)}
\]
follows directly from \eqref{eq:uniform-contour-heat}. All the
lower bounds cover complete edge segments. They are obtained by
subdividing each edge into dyadic intervals and evaluating
\[
 H_0(z)=\frac18(s-1)\pi^{-s/2}\Gamma(1+s/2)\zeta(s),
 \qquad s=\tfrac12+iz/2,
\]
as a complex ball on the entire segment. No numerical list of zeros
enters this computation.

The Newton integral is evaluated by certified complex quadrature,
using the exact logarithmic derivative
\[
 \frac{H_0'(z)}{H_0(z)}=\frac i2\left\{
 \frac1{s-1}-\frac12\log\pi
 +\frac12\psi(1+s/2)+\frac{\zeta'(s)}{\zeta(s)}\right\}.
\]
Its enclosures lie within $10^{-6}$ of the displayed integers;
the argument principle makes those integers exact. The real
endpoint values have opposite strict signs on the heat-coordinate
intervals
\[
 (28,29),\ (42,43),\ (50,51),\ (60,61),\
 (65,66),\ (75,76),\ (81,82),\ (86,87).
\]
The first rectangle uses just the first interval. For the second
rectangle the eight sign changes account for its entire degree,
counted with multiplicity. Thus every enclosed zero is simple, and
the distance between successive centers is greater than $4$.
These sign enclosures and the Newton integral also give a direct
verification of reality and simplicity on this particular rectangle.

For the second rectangle, $Z<91$ and $B_k<10^{240}$ for
$1\le k\le8$. Its leading coefficients satisfy
$D_k(0)\ge4^{k(k-1)}\ge1$. The Taylor remainder bound gives
\[
 |D_k(t)-D_k(0)|<\tfrac15
       \qquad (|t|\le10^{-257},\ 1\le k\le8).
\]
For real $t$ in this interval all the leading Hankel determinants
are positive. This implies that the eight roots are real and simple.
To verify that implication directly, a repeated root makes the
moment matrix singular. For distinct roots, a nonreal conjugate
pair admits a real polynomial of degree below eight taking the
values $i,-i$ there and zero at the other roots, by interpolation
and conjugation. Its moment square equals $-2$, contradicting
positive definiteness. For $T=20$ the constant enclosed degree is
one throughout $|t|\le10^{-8}$; the unique zero is therefore simple
throughout that complex disk and real at real heat times.

There is also an executed contour comparison for the ordinary
Hurwitz translation at height $20$. On the corresponding
$s$-rectangle, $3/8\le\Re s\le5/8$, $0\le\Im s\le20$,
so $|s|<21$. For $|u|\le1/4$, integration of each summand in $u$
gives
\[
 |\zeta(s,1+u)-\zeta(s)|
 \le 2|s|e^7|u|\sum_{n\ge1}n^{-11/8}
 <200000|u|.
\]
Here $|n+w|\ge3n/4$ and $|\arg(n+w)|\le1/(3n)$ on
the straight segment from $0$ to $u$; also
$\zeta(11/8)<1+8/3<4$. These estimates justify the convergent
difference series throughout this rectangle. The completion factor
$(s-1)\pi^{-s/2}\Gamma(1+s/2)/8$ has modulus less than $3$:
the Euler integral bounds the gamma modulus by
$\Gamma(1+\Re s/2)\le1$, the latter inequality following from
logarithmic convexity between $1$ and $2$. Hence the completed
translation differs from $H_0$ by less than $600000|u|$.
For $|u|\le10^{-12}$ this is smaller than the certified boundary
minimum. Rouch\'e's theorem gives exactly one simple enclosed zero
throughout that complex parameter disk. Averaging the same estimate
over any finite set of rotations gives the identical conclusion
for every rotation average in Theorem~\ref{thm:phase-orders}.

The computations use FLINT's rigorous complex-ball integration
\cite[section ``Integration'']{flint-acb-calc}. Their editable
program and interval receipts are in the workbench cited in the
closing note. The height-$20$ contour will be compared with an
executed probe-mass certificate in Section~\ref{sec:suzuki-probes}.
